\project{09}{CAN-FD from the controller out: bit timing and the first frame}{A voice}{Bit timing for both phases, the transceiver, sample point, the first frame on a wire}

\begin{keyfacts}
\item[Adds to the node] A voice. The node can put a frame on a wire, and can prove its own bit timing before a second node exists
\item[Peripherals] One bus controller instance, its message memory, its timestamp unit from chapter 3, and two pins once a transceiver is fitted
\item[Depends on] Chapter 4 for the kernel clock selector and the pin budget, chapter 3 for the timestamp at the start-of-frame bit
\item[Real or modelled] \textbf{Real}, with one caveat: everything up to the pins is real from the first evening, and the wire is real from the day the transceiver arrives. The controller's loopback modes are what make the gap survivable
\item[Difficulty] 4 of 5
\item[Effort] Three evenings before the transceiver, one after it
\item[Deliverable] Bit timing computed from the kernel clock with both phases checked for exact division, a message memory layout that is configured rather than assumed, a frame proven in internal loopback with no hardware at all, and the bit time confirmed by the controller's own timestamp counter
\end{keyfacts}

\subsection*{Why this chapter}

Eight chapters of sensing and acting produce a node that knows a great deal and
tells nobody. This chapter gives it a voice, and the first thing to say about
that voice is that it costs money: this board has no bus transceiver, and
without one the two pins are logic-level signals that no other node can hear.

That could have stalled the chapter until a part arrived. It does not, because
the controller has loopback modes: it can transmit a frame and receive its own
transmission without any external hardware at all. Nearly everything difficult
in this chapter, the bit timing arithmetic, the message memory layout, the
filters, the first frame, can be built and proven in that mode. What the
transceiver adds is the wire, the second node and the electrical reality, and
those arrive in chapter 10.

The second thing worth saying early is a correction. The sibling volume warns
that this part is not the family member most published material was written
for, and that warning is true at board level, clock level and clock-register
level. **At the bus peripheral it is false**, and that is good news: a
comparison of the vendor's own device headers shows the same licensed controller
core at the same addresses with the same register layout on both parts.
Material written for the popular sibling's bus peripheral is directly reusable
here. Three things are not, and chapter 4 has already dealt with the one that
matters.

\begin{note}{Buy the transceiver before starting this chapter}
Three facts that are worth more than a week. \textbf{A one megabit transceiver
still gives full sixty-four byte frames and the stronger checksum}, because the
frame format is invisible to the transceiver; what it costs is only the switch
to a faster data phase, which halves the demonstration rather than removing it.
\textbf{The part number does not tell you which you have}: one family carries no
format marker in its name at all and is rated five megabit, while another
separates one megabit parts from five megabit parts by a single letter suffix
inside the same datasheet. Read the suffix. And \textbf{a ready-made shield
exists} in the right form factor for this board, permissively licensed, carrying
a five megabit transceiver, from the same author as the bootloader project
chapter 19 discusses.
\end{note}

\subsection*{Prior art and what to reuse}

\begin{table}[H]\centering\small
\begin{tabular}{p{34mm}p{46mm}p{38mm}p{20mm}}
\toprule
Source & What it gives & What it does not & Licence \\
\midrule
The silicon vendor's device headers for this part and for its sibling & The correction this chapter opens with, in a form a reader can repeat in ten minutes: the same controller core, the same two base addresses, the same calibration unit and message memory base, and identical bit-timing field positions on both parts & They do not name the licensed core's own documentation, which is where the field semantics are defined & Apache-2.0 \\
The vendor's application note on this peripheral & The introduction to the peripheral for this microcontroller family & \textbf{Indexed only.} It could not be fetched during the research for this volume, so it is cited by number and title and nothing is quoted from it & vendor document \\
A ready-made transceiver shield for these boards & A five megabit transceiver in the right form factor, with a published design, from the same author as the bootloader project in chapter 19 & It still has to be bought or built, and it occupies the header the sensor shield uses & MIT \\
A bit timing calculator & The arithmetic for both phases, laid out, with the sample point and the tolerance window & \textbf{Copyleft over a network}, so use the hosted page and do not vendor it. The arithmetic in step 3 is written out here so that the node does not depend on it at all & AGPL-3.0 \\
The upstream operating system's board files & The cleanest published statement that these boards carry no transceiver, and, for the sibling board, a working pin pair for the peripheral & The file for this board does not mention the peripheral at all, which is itself the evidence & Apache-2.0 \\
\bottomrule
\end{tabular}
\caption{Prior art for chapter 9. The fourth row is the licence decision of the chapter: a calculator under a network copyleft licence is fine to use through a browser and is not something to copy into a repository, so the arithmetic is written out instead.}
\end{table}

\subsection*{What the node gains}

Before this chapter the node is silent. After it, the node can frame a message,
put it on its own transmit pin, receive it, and state its bit timing in both
phases with the arithmetic shown and the result confirmed by the controller's
own timestamp counter. It cannot yet reach another machine, because that needs
the transceiver and the wire, which is chapter 10.

\subsection*{Parts from the inventory}

\begin{table}[H]\centering\small
\begin{tabular}{p{40mm}p{66mm}p{40mm}}
\toprule
Part & Role & Interface \\
\midrule
NUCLEO-H7A3ZI-Q & The controller, its message memory and its timestamp unit & Micro USB to the host \\
\textbf{To be bought:} a CAN-FD transceiver & The only electrical gap between this node and a bus. Read the suffix & Two logic pins, and the bus pair \\
\textbf{To be bought:} two 120 ohm terminators & One at each end of the bus, not one per node & Across the bus pair \\
Two jumper wires & The external loopback in step 6, before a second node exists & Transmit pin to receive pin \\
\bottomrule
\end{tabular}
\caption{Inventory items used in chapter 9. Two of the four have to be bought, and the chapter is arranged so that three of its seven steps can be completed before they arrive.}
\end{table}

\subsection*{System architecture}

\diagram{j09_arch}{The controller, its message memory, and the three ways a
frame can travel. Internal loopback needs no pins at all; external loopback
needs two pins and one wire; the bus needs the transceiver and a second node.
The chapter is built in that order, which is also the order the parts arrive in.}

\subsection*{Peripheral configuration}

\begin{table}[H]\centering\small
\begin{tabular}{p{22mm}p{26mm}p{24mm}p{30mm}p{30mm}}
\toprule
Peripheral & Mode & Clock source & Pins and function & Interrupt and transfers \\
\midrule
Bus controller, first instance & Classic, then flexible-data without and with the rate switch & Kernel clock, selected by the register named in chapter 4 & Two pins on port D, confirmed against the board manual & Receive and transmit interrupts, plus the error and bus-off interrupt \\
Its message memory & Configured: filters, receive buffers, transmit buffers & n/a & None & None. It is memory, and it must be laid out before use \\
Its timestamp unit & Counting bit times & The bus bit time & None & None. Read in step 7 \\
GPIO & Alternate function, and open drain on the transmit pin where the transceiver needs it & Peripheral bus & Two pins & None \\
\bottomrule
\end{tabular}
\caption{Peripheral configuration for chapter 9. The second row is the one that produces the classic symptom: a controller whose message memory has not been laid out accepts configuration, reports no error, and never transmits.}
\end{table}

\subsection*{Wiring}

\diagram{j09_wiring}{The bus as it will be, and the two things to get right when
it is built. Termination goes at the two ends of the bus and not on every node,
and the transceiver's suffix decides whether the data phase can run faster than
the arbitration phase at all.}

\subsection*{Memory and timing budget}

\begin{table}[H]\centering\small
\begin{tabular}{p{44mm}p{28mm}p{28mm}p{30mm}}
\toprule
Quantity & Budget & Measured & Margin \\
\midrule
Flash, this chapter & 9 kB & not measured & not measured \\
Message memory used & under 2 kB & computed from the layout & n/a \\
Nominal bit rate & 500 kbit/s & confirmed in step 7 & n/a \\
Data bit rate & 2 Mbit/s & confirmed in step 7 & n/a \\
Sample point, nominal phase & 80 per cent & computed & n/a \\
Sample point, data phase & 75 per cent & computed & n/a \\
Frame duration, 64 bytes at both rates & computed in step 7 & not measured & not measured \\
\bottomrule
\end{tabular}
\caption{The budget for chapter 9. The two rate rows say confirmed rather than measured, because the confirmation comes from the controller's own timestamp counter counting bit times rather than from an instrument on the wire, and the chapter is explicit about the difference.}
\end{table}

\subsection*{Firmware design (UML)}

\diagram{j09_uml}{The bring-up order, and the three checks that stop the chapter
rather than letting it continue with a fault. Each check is cheap, each catches
a defect that otherwise appears much later as a silent failure to transmit, and
the last two of the five stages need hardware that has to be bought.}

Two design decisions carry into the rest of the volume.

\textbf{The bit timing is computed, not copied.} A table of register values
found online encodes somebody else's kernel clock. The node computes its
timing from the kernel clock it reads back at boot, refuses to start if either
phase does not divide exactly, and prints both results with their sample points.

\textbf{Transmission never blocks the control loop.} A frame is handed to the
controller and the controller sends it when the bus allows. The period handler
from chapter 2 writes into a transmit buffer and returns; if the buffer is full
it counts that and returns anyway. A control loop that waits for bus arbitration
has given the bus control over its deadline.

\subsection*{Data flow (ASCII)}

\begin{asciiart}
  kernel clock, read back from the registers at boot
        |
        v
  +---------------------------------------------------------------+
  | nominal phase:  prescaler -> time quanta -> segments           |
  |                 500 kbit/s, sample point 80 per cent            |
  | data phase:     its own prescaler and segments                  |
  |                 2 Mbit/s, sample point 75 per cent              |
  |                 plus transmitter delay compensation             |
  +---------------------------------------------------------------+
        |  exact division, or the node refuses to start
        v
  +---------------------------------------------------------------+
  | message memory, laid out before anything is sent:              |
  |   standard filters | extended filters | receive FIFO 0 |        |
  |   receive FIFO 1   | receive buffers  | transmit event |        |
  |   transmit buffers                                              |
  +---------------------------------------------------------------+
        |
        v
  internal loopback  --> no pins, no hardware, and most of the chapter
  external loopback  --> two pins and one wire
  the bus            --> transceiver, termination, a second node (chapter 10)
        |
        v
  every received frame carries a timestamp captured at its start-of-frame bit
  (chapter 3), which is what chapter 12 needs
\end{asciiart}

\subsection*{Repository layout}

\begin{plaincode}
joint-node/
  board/  nucleo_h7a3zi_q.yaml        # the two bus pins, evidence now confirmed
  doc/    bit-timing.md               # + this chapter: both phases, computed
  src/
    bus/
      fdcan.c     fdcan.h             # + this chapter: bring-up and send
      bittiming.c bittiming.h         # + this chapter: the arithmetic
      msgram.c    msgram.h            # + this chapter: the memory layout
      loopback.c  loopback.h          # + this chapter: the three modes
    sense/ act/ estimate/ control/ time/ node/ bsp/
    mw/  safety/  update/
  test/
    test_bittiming.c                  # + the arithmetic, on the host
    test_msgram.c                     # + the layout, on the host
  host/  tools/  README.md
\end{plaincode}

\subsection*{Steps}

\begin{steps}
\item \textbf{Select and read back the kernel clock.} The peripheral's clock
does not come from the peripheral bus: it has its own selector, in the register
chapter 4 identified as the one whose name differs from the sibling part's. Pick
a source that divides cleanly into both bit rates, and read the result back
rather than trusting the selection.
\begin{ccode}
/* fdcan.c: the kernel clock is the foundation of every number below it. */
uint32_t fdcan_kernel_hz(void)
{
    switch (rcc_fdcan_source()) {                /* read back, do not assume */
        case FDCAN_SRC_HSE:  return board_clock_hz(BOARD_CLK_HSE);
        case FDCAN_SRC_PLL1Q: return board_clock_hz(BOARD_CLK_PLL1Q);
        case FDCAN_SRC_PLL2Q: return board_clock_hz(BOARD_CLK_PLL2Q);
        default: fail("fdcan: no kernel clock selected");
    }
}
\end{ccode}

\item \textbf{Compute the bit timing, and refuse anything inexact.} A bit is
divided into time quanta; the quantum is the kernel clock divided by a
prescaler; the bit is a synchronisation quantum plus two segments; and the
sample point is where the second segment begins. Both phases have their own
prescaler and their own segments.
\begin{ccode}
/* bittiming.c: the whole arithmetic, in one place, with no table of magic.
   The limits are an argument because the two phases do not share them. */
bool bt_compute(uint32_t kernel_hz, uint32_t bitrate, float want_sp,
                const bt_limits_t *lim, bt_t *out)
{
    uint32_t want = (uint32_t) (want_sp * 1000.0f + 0.5f);   /* per mille */
    uint32_t max_tq = 1u + lim->seg1_max + lim->seg2_max;

    for (uint32_t pre = lim->prescaler_min; pre <= lim->prescaler_max; pre++) {
        if (kernel_hz % pre) continue;                 /* inexact prescaler */
        uint32_t tq_hz = kernel_hz / pre;
        if (tq_hz % bitrate) continue;                 /* inexact bit time  */
        uint32_t tq_per_bit = tq_hz / bitrate;
        if (tq_per_bit < lim->min_tq || tq_per_bit > max_tq) continue;

        uint32_t before = (want * tq_per_bit + 500u) / 1000u;
        if (before == 0u || before > tq_per_bit) continue;
        uint32_t seg1 = before - 1u;                   /* 1 tq is the sync  */
        uint32_t seg2 = tq_per_bit - seg1 - 1u;
        if (seg1 < lim->seg1_min || seg1 > lim->seg1_max) continue;
        if (seg2 < lim->seg2_min || seg2 > lim->seg2_max) continue;

        out->prescaler = pre; out->seg1 = seg1; out->seg2 = seg2;
        out->tq_per_bit = tq_per_bit;
        out->sjw = min_u32(seg2, lim->sjw_max);
        out->sample_point_permille =
            ((1u + seg1) * 1000u + tq_per_bit / 2u) / tq_per_bit;
        return true;
    }
    return false;                                      /* caller must fail   */
}
\end{ccode}
Print both results at boot, with their achieved sample points, and stop the node
if either phase could not be solved exactly. A bit rate that is one part in a
thousand out works between two nodes that are wrong in the same direction and
fails against anything else.

Two things in that function are easy to get wrong and both cost a working bus.
The first is that \textbf{the limits are an argument rather than a set of
constants}, because the registers that hold these numbers are much narrower in
the data phase: the prescaler there runs to 32 where the nominal phase allows
512, and the second segment to 16 where the nominal phase allows 128. One set of
constants cannot be right for both, and the set that suits the nominal phase
silently accepts data timings the hardware cannot hold.

The second is that \textbf{the decision is integer}. The obvious way to place
the sample point is to round a fraction of the bit, and the obvious way to check
the result is to compute the same fraction in whatever language the host test is
written in. Those two roundings do not have to agree: one common pair of
conventions rounds a half to the nearest even number and the other rounds it
away from zero, and 75 per cent of 30 quanta is exactly 22.5. That disagreement
produces one wrong row in a table of correct ones, which is the hardest kind to
find.

\diagram{j09_data}{One bit in each phase, divided into time quanta, with the
sample point marked. The two phases share a kernel clock and nothing else: each
has its own prescaler and its own segment counts, which is why the arithmetic is
done twice and why both results have to divide exactly. The arithmetic is
written out here because the calculator that does it well is under a network
copyleft licence, so it is used through a browser and never vendored.}

\item \textbf{Lay out the message memory before configuring anything else.} On
this controller the filters, the receive buffers and the transmit buffers all
live in one block of memory whose layout the firmware chooses. A controller
whose layout has not been written accepts every other configuration, reports no
error, and never transmits, which is the single most common way this peripheral
appears broken.
\begin{ccode}
/* msgram.c: an explicit layout, computed once and asserted against the size. */
static const msgram_layout_t layout = {
    .std_filters = 8,      .ext_filters = 4,
    .rx_fifo0    = 16,     .rx_fifo1    = 8,
    .rx_buffers  = 0,      .tx_event    = 8,
    .tx_buffers  = 8,      .element_bytes = 72,   /* 64 data + header  */
};
/* The build fails if this layout does not fit the part's message memory. */
_Static_assert(MSGRAM_BYTES(layout) <= MSGRAM_SIZE_BYTES, "message memory");
\end{ccode}

\item \textbf{Send a frame to yourself, with no hardware at all.} Internal
loopback connects the controller's transmitter to its own receiver inside the
peripheral. The pins are not involved, no transceiver is needed and no second
node exists. Almost everything in this chapter can be proven here.
\begin{ccode}
fdcan_mode(FDCAN_MODE_INTERNAL_LOOPBACK);
fdcan_send(&(frame_t){ .id = 0x123, .len = 8, .fd = false, .brs = false, ... });
frame_t got;
assert(fdcan_receive(&got, 10 /* ms */));
assert(got.id == 0x123 && got.len == 8 && memcmp(got.data, sent.data, 8) == 0);
\end{ccode}
Run the same test three times: a classic frame, a flexible-data frame without
the rate switch, and a flexible-data frame with it. The third exercises the data
phase timing computed in step 2, which is the part a classic-only setup never
touches.

\item \textbf{Turn on the transmitter delay compensation before the data phase
goes fast.} At data rates above roughly one megabit the round trip through the
transceiver becomes a significant fraction of a bit time, and the controller has
to compensate for it or it will sample its own transmission in the wrong place.
The offset comes from the transceiver's loop delay, which is in its datasheet,
and the controller can also measure it.
\begin{plaincode}
transceiver loop delay   from the datasheet, in nanoseconds
                         converted to time quanta at the data-phase quantum
compensation offset      set from that, or measured by the controller itself
when it matters          above about 1 Mbit/s in the data phase
symptom when wrong       frames that pass in loopback and fail on a real bus,
                         with errors that increase with the data rate
\end{plaincode}

\item \textbf{Prove it on the pins with external loopback and one wire.} The
external mode drives the transmit pin and listens on the receive pin, so a
single jumper wire between them proves the pin configuration, the alternate
functions and the direction, which internal loopback cannot.

\item \textbf{Confirm the bit time with the controller's own timestamp
counter.} This is the chapter's measurement, and it needs no instrument. The
timestamp unit can be clocked from the bus bit time rather than from a timer.
Configure it that way, send frames of known length, and read the difference
between their captured timestamps: the result is the frame's duration in bit
times, which confirms both the bit timing and the frame format.
\begin{shellcode}
python host/frame_timing.py --lengths 0,8,16,32,64 --fd --brs
# len  measured bit times  expected  note
#   0                  67        67  arbitration only
#   8                 132       132  data phase at 2 Mbit/s
#  16                 197       197
#  32                 327       327
#  64                 587       587  and the timestamp is taken at the SOF bit
\end{shellcode}
Two honest notes. The counter counts bit times, so this confirms the ratio
between the phases and the frame structure rather than the absolute rate in
bits per second; the absolute rate rests on the kernel clock, which chapter 4
reads back from the registers. And the timestamps are captured at the
start-of-frame bit, in hardware, which is the property chapter 3 established and
chapter 12 depends on.

\item \textbf{Fit the transceiver, terminate the bus, and send one frame to
nothing.} With the transceiver fitted and a single terminator, the node is alone
on a bus and every frame it sends goes unacknowledged. That is not a failure to
debug: a frame with no other node to acknowledge it is retransmitted, the error
counter rises, and the controller eventually becomes error-passive. Watching
that happen deliberately is the best possible preparation for chapter 12.
\end{steps}

\subsection*{Build, flash and debug}

\diagram{j09_timing}{One flexible-data frame with the rate switch set, drawn
against one classic frame of the same payload. The arbitration phase runs at the
nominal rate in both; only the data phase changes speed, and only between the
two points marked. The timestamp the controller captures is at the start of the
frame, before any of this.}

\begin{shellcode}
cmake --build build -j && probe-rs run --chip STM32H7A3ZITx build/firmware.elf
python host/frame_timing.py --lengths 0,8,16,32,64 --fd --brs
\end{shellcode}

\begin{note}{When the controller accepts everything and transmits nothing}
In order of likelihood: the message memory was never laid out, so the transmit
buffer the firmware writes into is not where the controller looks; the
controller was left in its initialisation state, because leaving it requires an
explicit step that is easy to omit; the kernel clock selector was never set, so
the peripheral has no clock at all and its registers still read back the values
that were written; or the node is alone on a terminated bus and is
retransmitting a frame nobody will acknowledge, which looks identical from the
firmware's side until the error counters are read. The last one is step 8 and is
worth doing on purpose.
\end{note}

\subsection*{Verification and acceptance criteria}

\begin{itemize}
\item The kernel clock is read back from the registers and printed, and the node
refuses to start if no source is selected.
\item Both phases' timing is computed from that clock, both divide exactly, and
both achieved sample points are printed. A deliberately impossible bit rate
stops the node with a named error rather than being approximated.
\item The bit timing arithmetic has host tests over a range of kernel clocks and
bit rates, including the cases that must fail.
\item The message memory layout is asserted against the part's memory size at
compile time, and the assertion is proven by deliberately enlarging the layout.
\item A classic frame, a flexible-data frame and a flexible-data frame with the
rate switch all complete in internal loopback with the payload compared byte by
byte.
\item External loopback with one jumper wire passes the same three tests, which
proves the pins as well as the controller.
\item The measured frame durations in bit times match the computed ones for
every length in the sweep.
\item The period handler's transmit path never blocks: a full transmit buffer
increments a counter and the handler returns within its budget, proven by
filling the buffer deliberately.
\end{itemize}

\subsection*{Variants}

\begin{table}[H]\centering\small
\begin{tabular}{p{24mm}p{26mm}p{44mm}p{22mm}p{20mm}}
\toprule
Axis & Variant & What changes & Cost & Built in full in \\
\midrule
Frame format & Classic, 8 bytes & Works with every node ever built, and the checksum is the weaker one & Payload, and the checksum & Here \\
Frame format & Flexible-data, no rate switch & Sixty-four byte payload and the stronger checksum, at the arbitration rate throughout & Nothing, on a transceiver of any rating & Here \\
Frame format & Flexible-data with the rate switch & The data phase runs faster than arbitration, which is the whole point of the format & A transceiver rated for the data rate, and delay compensation & Here \\
Transceiver & Rated 1 Mbit & Full sixty-four byte frames and the stronger checksum, with no rate switch & Half the demonstration, and it is the cheaper part & Named, and the chapter says which was used \\
Transceiver & Rated 5 Mbit & Everything, including the rate switch & Read the suffix & Here \\
Loopback & Internal & No pins, no hardware, and most of this chapter & It proves nothing about the pins & Here \\
Loopback & External, one wire & Proves the pins, the alternate functions and the direction & One wire & Here \\
Delivery & Polled from the loop & Simplest, and the loop then waits for the bus & The deadline & Nowhere. It is the thing the design rule forbids \\
Delivery & Interrupt into a buffer & The baseline: the loop writes and returns & A buffer, and a full-buffer policy & Here \\
Delivery & Transfers into message memory & The controller's own memory is filled without the processor & Complexity, for a gain this node does not need at its frame rate & Chapter 13 \\
\bottomrule
\end{tabular}
\caption{Variants for chapter 9. The three frame-format rows are worth building in that order even with a fast transceiver in hand, because each one exercises a different part of the timing arithmetic.}
\end{table}

\subsection*{Pitfalls}

\begin{itemize}
\item Copying register values from an example. They encode somebody else's
kernel clock, and the symptom is a node that works against its author's bench
and nothing else.
\item Leaving the message memory unconfigured. Everything else succeeds and
nothing is ever transmitted.
\item Forgetting to leave the controller's initialisation state. The registers
read back correctly and the bus stays silent.
\item Terminating every node. Termination belongs at the two ends of the bus,
and a third terminator loads the bus enough to stop it working at speed.
\item Buying a transceiver by family name. One family's name carries no format
marker at all, and another separates one megabit from five megabit parts by a
single letter in the same datasheet.
\item Running the data phase above about one megabit without delay
compensation. It passes in loopback and fails on a real bus, and the failure
rate rises with the data rate.
\item Waiting for transmission in the control loop. The bus then owns the
node's deadline.
\item Reading a bit rate confirmed by the controller's own counter as an
absolute measurement. It confirms the structure and the ratio; the absolute rate
rests on the kernel clock.
\end{itemize}

\subsection*{Best practices applied}

\begin{itemize}
\item Every number is computed from a clock the firmware read back, and an
inexact result stops the node rather than being rounded.
\item The chapter is ordered so that the parts that have to be bought are needed
last, and three of its steps complete before they arrive.
\item The peripheral's own facilities are used as the instrument: the timestamp
counter clocked from the bit time confirms the frame structure with nothing
attached to the board.
\item A licence decision is made explicitly rather than by habit: a useful
calculator under a network copyleft licence is used through a browser, and the
arithmetic is written out so the node does not depend on it.
\item The correction to the volume's own part-trap warning is stated where it
applies, with evidence a reader can repeat.
\item A deliberate failure, a node alone on a bus, is performed on purpose
because the error behaviour it produces is chapter 12's subject.
\end{itemize}

\subsection*{Stretch goals}

\begin{itemize}
\item Have the controller measure the transceiver's loop delay itself and
compare it with the datasheet figure. That is a real measurement of a real part,
and it is one of the few in this volume that needs no additional instrument.
\item Sweep the sample point across its usable range at a fixed bit rate and
find where communication starts to fail on a real bus with a long stub. That
figure is rarely published and is easy to produce once two nodes exist.
\item Compute the bit timing for every kernel clock the part can produce and
publish the table of which bit rates are exactly achievable. It is a small
piece of genuinely useful reference material for this family.
\item Build the transceiver shield from the published design rather than buying
a module, and compare the two for propagation delay.
\end{itemize}

\subsection*{Roadmap and next steps}

Chapter 10 builds the other end: the motion master on a Raspberry Pi, the
kernel's own socket layer, the command line tools, and the adapter decision that
turns out to matter more than it looks.

The published progression from here is the bus association's own documents and
its conference proceedings, which are the primary literature for this subject
and are largely free to read. The data link standard's current edition is the
normative statement of the frame format and is paywalled; its 2015 predecessor
carried a different subtitle, so the edition matters when citing it. For the
peripheral itself, the vendor's application note is the right next document and
is an existence claim only in this volume, because it could not be fetched
during the research.

\subsection*{Portfolio evidence}

\begin{itemize}
\item The computed bit timing for both phases, with the achieved sample points
and the exact-division check, printed at boot.
\item The frame duration table: measured bit times against computed ones for
five payload lengths, produced with no instrument attached to the board.
\item The three loopback tests, which demonstrate that most of a bus bring-up
can be completed and proven before the hardware arrives.
\item The header comparison from the chapter's opening, which corrects a warning
this volume itself carries and shows the evidence for the correction.
\end{itemize}

\subsection*{Sources}

Normative references:
\begin{itemize}
\item ISO 11898-1:2024, the data link layer and physical coding sublayer, third
edition, \pubdate{May 2024}, which supersedes the 2015 edition whose subtitle differed.
Paywalled; cited by number, title and edition.
\item Reference manual RM0455, for the controller's registers, its message
memory and its timestamp unit.
\item The vendor's application note on this peripheral for this microcontroller
family, cited as an existence claim only.
\item The transceiver's datasheet, for the loop delay used in step 5 and for the
suffix that decides the maximum data rate.
\end{itemize}

Reusable implementations:
\begin{itemize}
\item The vendor's device headers, Apache-2.0, which settle that the controller
core is the same on this part and on its better-known sibling.\newline
\url{https://github.com/STMicroelectronics/cmsis-device-h7}
\item A ready-made transceiver shield for these boards, MIT.\newline
\url{https://github.com/feaser/canfdshield}
\item The upstream operating system's board files, Apache-2.0, which state
plainly that these boards carry no transceiver.\newline
\url{https://docs.zephyrproject.org/latest/boards/st/nucleo_h7a3zi_q/doc/index.html}
\end{itemize}
